\begingroup
\singlespacing
\fontsize{12}{14.4}\selectfont
\RaggedRight
\addcontentsline{toc}{section}{Список литературы}
\begin{thebibliography}{99}

\bibitem{sze2017}
Sze V., Chen Y.-H., Yang T.-J., Emer J.~S.
Efficient Processing of Deep Neural Networks: A Tutorial and Survey //
Proceedings of the IEEE. 2017. Vol.~105, no.~12. P.~2295--2329.
\href{https://arxiv.org/abs/1703.09039}{arXiv:1703.09039}.

\bibitem{bush1931}
Bush V.
The Differential Analyzer. A New Machine for Solving Differential Equations //
Journal of the Franklin Institute. 1931. Vol.~212, no.~4. P.~447--488.
\href{https://fab.cba.mit.edu/classes/862.22/notes/computation/Bush-1931.pdf}{Полный текст статьи}.

\bibitem{mcculloch1943}
McCulloch W.~S., Pitts W.
A Logical Calculus of the Ideas Immanent in Nervous Activity //
Bulletin of Mathematical Biophysics. 1943. Vol.~5. P.~115--133.
\href{https://doi.org/10.1007/BF02478259}{DOI: 10.1007/BF02478259}.

\bibitem{rosenblatt1958}
Rosenblatt F.
The Perceptron: A Probabilistic Model for Information Storage and
Organization in the Brain // Psychological Review. 1958.
Vol.~65, no.~6. P.~386--408.
\href{https://www.cs.cmu.edu/~epxing/Class/10715/reading/Rosenblatt.perceptron.pdf}{Полный текст статьи}.

\bibitem{widrow1960memistor}
Widrow B.
An Adaptive ``Adaline'' Neuron Using Chemical ``Memistors''.
Stanford Electronics Laboratories, Technical Report 1553-2. October 1960.
\href{https://isl.stanford.edu/~widrow/papers/t1960anadaptive.pdf}{Текст отчёта у автора}.

\bibitem{widrow1995}
Widrow B., Lehr M.~A.
Perceptrons, Adalines, and Backpropagation //
The Handbook of Brain Theory and Neural Networks / Ed. M.~A. Arbib.
MIT Press, 1995. P.~719--724.
\href{https://www-isl.stanford.edu/~widrow/papers/bc1995perceptronsadalines.pdf}{Текст главы у автора}.

\bibitem{landauer1961}
Landauer R.
Irreversibility and Heat Generation in the Computing Process //
IBM Journal of Research and Development. 1961. Vol.~5, no.~3. P.~183--191.
\href{https://doi.org/10.1147/rd.53.0183}{DOI: 10.1147/rd.53.0183}.

\bibitem{bennett1973}
Bennett C.~H.
Logical Reversibility of Computation //
IBM Journal of Research and Development. 1973. Vol.~17, no.~6. P.~525--532.
\href{https://doi.org/10.1147/rd.176.0525}{DOI: 10.1147/rd.176.0525}.

\bibitem{wanlass1963}
Wanlass F.~M., Sah C.~T.
Nanowatt Logic Using Field-Effect Metal-Oxide Semiconductor Triodes //
IEEE International Solid-State Circuits Conference. Digest of Technical Papers.
1963. P.~32--33.
\href{https://doi.org/10.1109/ISSCC.1963.1157450}{DOI: 10.1109/ISSCC.1963.1157450}.

\bibitem{rumelhart1986}
Rumelhart D.~E., Hinton G.~E., Williams R.~J.
Learning Representations by Back-Propagating Errors //
Nature. 1986. Vol.~323. P.~533--536.
\href{https://doi.org/10.1038/323533a0}{DOI: 10.1038/323533a0}.

\bibitem{lecun1998}
LeCun Y., Bottou L., Bengio Y., Haffner P.
Gradient-Based Learning Applied to Document Recognition //
Proceedings of the IEEE. 1998. Vol.~86, no.~11. P.~2278--2324.
\href{https://bottou.org/papers/lecun-98h}{Страница публикации у соавтора}.

\bibitem{wilkes1965}
Wilkes M.~V.
Slave Memories and Dynamic Storage Allocation //
IEEE Transactions on Electronic Computers. 1965. Vol.~EC-14, no.~2. P.~270--271.
\href{https://www.cl.cam.ac.uk/teaching/2425/P56/files/wilkes1965cache.pdf}{Полный текст статьи}.

\bibitem{amdahl1967}
Amdahl G.~M.
Validity of the Single Processor Approach to Achieving Large Scale
Computing Capabilities // AFIPS Spring Joint Computer Conference.
1967. P.~483--485.
\href{https://www.cs.cmu.edu/~18742/papers/Amdahl1967.pdf}{Полный текст статьи}.

\bibitem{kung1982}
Kung H.~T. Why Systolic Architectures? //
Computer. 1982. Vol.~15, no.~1. P.~37--46.
\href{https://doi.org/10.1109/MC.1982.1653825}{DOI: 10.1109/MC.1982.1653825}.

\bibitem{hong1981}
Hong J.-W., Kung H.~T.
I/O Complexity: The Red-Blue Pebble Game //
Proceedings of the 13th Annual ACM Symposium on Theory of Computing.
1981. P.~326--333.
\href{https://doi.org/10.1145/800076.802486}{DOI: 10.1145/800076.802486}.

\bibitem{mead1990}
Mead C.
Neuromorphic Electronic Systems //
Proceedings of the IEEE. 1990. Vol.~78, no.~10. P.~1629--1636.
\href{https://authors.library.caltech.edu/records/j6gtc-ptx47}{Статья в архиве Caltech}.

\bibitem{goldberg1991}
Goldberg D.
What Every Computer Scientist Should Know About Floating-Point Arithmetic //
ACM Computing Surveys. 1991. Vol.~23, no.~1. P.~5--48.
\href{https://doi.org/10.1145/103162.103163}{DOI: 10.1145/103162.103163}.

\bibitem{dennard1974}
Dennard R.~H., Gaensslen F.~H., Yu H.-N. et al.
Design of Ion-Implanted MOSFET's with Very Small Physical Dimensions //
IEEE Journal of Solid-State Circuits. 1974. Vol.~9, no.~5. P.~256--268.
\href{https://doi.org/10.1109/JSSC.1974.1050511}{DOI: 10.1109/JSSC.1974.1050511}.

\bibitem{dark2012}
Esmaeilzadeh H., Blem E., St.~Amant R. et al.
Power Limitations and Dark Silicon Challenge the Future of Multicore //
ACM Transactions on Computer Systems. 2012. Vol.~30, no.~3. Article~11.
\href{https://doi.org/10.1145/2324876.2324879}{DOI: 10.1145/2324876.2324879}.

\bibitem{wulf1995}
Wulf W.~A., McKee S.~A.
Hitting the Memory Wall: Implications of the Obvious //
ACM SIGARCH Computer Architecture News. 1995. Vol.~23, no.~1. P.~20--24.
\href{https://doi.org/10.1145/216585.216588}{DOI: 10.1145/216585.216588}.

\bibitem{krizhevsky2012}
Krizhevsky A., Sutskever I., Hinton G.~E.
ImageNet Classification with Deep Convolutional Neural Networks //
Advances in Neural Information Processing Systems. 2012. Vol.~25.
\href{https://proceedings.neurips.cc/paper_files/paper/2012/hash/c399862d3b9d6b76c8436e924a68c45b-Abstract.html}{Страница статьи в материалах конференции}.

\bibitem{chen2014}
Chen T., Du Z., Sun N. et al.
DianNao: A Small-Footprint High-Throughput Accelerator for Ubiquitous
Machine-Learning // Proceedings of ASPLOS. 2014.
\href{https://users.cs.duke.edu/~lkw34/papers/diannao-asplos2014.pdf}{Полный текст статьи}.

\bibitem{eyeriss2016}
Chen Y.-H., Emer J., Sze V.
Eyeriss: A Spatial Architecture for Energy-Efficient Dataflow for
Convolutional Neural Networks // Proceedings of ISCA. 2016. P.~367--379.
\href{https://eyeriss.mit.edu/}{Страница проекта и публикаций MIT}.

\bibitem{jouppi2017}
Jouppi N.~P., Young C., Patil N. et al.
In-Datacenter Performance Analysis of a Tensor Processing Unit //
Proceedings of ISCA. 2017.
\href{https://arxiv.org/abs/1704.04760}{arXiv:1704.04760}.

\bibitem{horowitz2014}
Horowitz M.
Computing's Energy Problem (and What We Can Do About It) //
IEEE International Solid-State Circuits Conference Digest of Technical Papers.
2014. P.~10--14.
\href{https://doi.org/10.1109/ISSCC.2014.6757323}{DOI: 10.1109/ISSCC.2014.6757323}.

\bibitem{roofline2009}
Williams S., Waterman A., Patterson D.
Roofline: An Insightful Visual Performance Model for Multicore Architectures //
Communications of the ACM. 2009. Vol.~52, no.~4. P.~65--76.
\href{https://doi.org/10.1145/1498765.1498785}{DOI: 10.1145/1498765.1498785}.

\bibitem{timeloop2019}
Parashar A., Raina P., Shao Y.~S. et al.
Timeloop: A Systematic Approach to DNN Accelerator Evaluation //
IEEE International Symposium on Performance Analysis of Systems and Software.
2019.
\href{https://research.nvidia.com/publication/2019-03_timeloop-systematic-approach-dnn-accelerator-evaluation}{Страница публикации NVIDIA Research}.

\bibitem{accelergy2019}
Wu Y.~N., Emer J.~S., Sze V.
Accelergy: An Architecture-Level Energy Estimation Methodology for
Accelerator Designs // International Conference on Computer-Aided Design.
2019.
\href{https://d1qx31qr3h6wln.cloudfront.net/publications/ICCAD_2019_Accelergy.pdf}{Полный текст статьи}.

\bibitem{obd1989}
LeCun Y., Denker J.~S., Solla S.~A.
Optimal Brain Damage // Advances in Neural Information Processing Systems.
Vol.~2. Материалы конференции 1989 года.
\href{https://proceedings.neurips.cc/paper_files/paper/1989/file/6c9882bbac1c7093bd25041881277658-Paper.pdf}{Полный текст статьи}.

\bibitem{gong2014}
Gong Y., Liu L., Yang M., Bourdev L.
Compressing Deep Convolutional Networks Using Vector Quantization.
2014. \href{https://arxiv.org/abs/1412.6115}{arXiv:1412.6115}.

\bibitem{han2016}
Han S., Mao H., Dally W.~J.
Deep Compression: Compressing Deep Neural Networks with Pruning,
Trained Quantization and Huffman Coding // ICLR. 2016.
Электронная публикация: 2015.
\href{https://arxiv.org/abs/1510.00149}{arXiv:1510.00149}.

\bibitem{hinton2015distillation}
Hinton G., Vinyals O., Dean J.
Distilling the Knowledge in a Neural Network. 2015.
\href{https://arxiv.org/abs/1503.02531}{arXiv:1503.02531}.

\bibitem{howard2017mobilenets}
Howard A.~G., Zhu M., Chen B. et al.
MobileNets: Efficient Convolutional Neural Networks for Mobile Vision
Applications. 2017.
\href{https://arxiv.org/abs/1704.04861}{arXiv:1704.04861}.

\bibitem{sung2016}
Sung W., Shin S., Hwang K.
Resiliency of Deep Neural Networks under Quantization.
Электронная публикация: 2015, версия~3: 2016.
\href{https://arxiv.org/abs/1511.06488v3}{arXiv:1511.06488v3}.

\bibitem{bnn2016}
Courbariaux M., Hubara I., Soudry D. et al.
Binarized Neural Networks: Training Deep Neural Networks with Weights
and Activations Constrained to +1 or $-1$.
2016. \href{https://arxiv.org/abs/1602.02830v3}{arXiv:1602.02830v3}.

\bibitem{jacob2018}
Jacob B., Kligys S., Chen B. et al.
Quantization and Training of Neural Networks for Efficient
Integer-Arithmetic-Only Inference // CVPR. 2018. Электронная публикация: 2017.
\href{https://arxiv.org/abs/1712.05877}{arXiv:1712.05877}.

\bibitem{mixed2018}
Micikevicius P., Narang S., Alben J. et al.
Mixed Precision Training // ICLR. 2018. Электронная публикация: 2017.
\href{https://arxiv.org/abs/1710.03740}{arXiv:1710.03740}.

\bibitem{hawq2019}
Dong Z., Yao Z., Gholami A., Mahoney M., Keutzer K.
HAWQ: Hessian AWare Quantization of Neural Networks with Mixed-Precision //
ICCV. 2019.
\href{https://arxiv.org/abs/1905.03696}{arXiv:1905.03696}.

\bibitem{haq2019}
Wang K., Liu Z., Lin Y., Lin J., Han S.
HAQ: Hardware-Aware Automated Quantization with Mixed Precision //
CVPR. 2019. P.~8612--8620.
\href{https://openaccess.thecvf.com/content_CVPR_2019/html/Wang_HAQ_Hardware-Aware_Automated_Quantization_With_Mixed_Precision_CVPR_2019_paper.html}{Страница статьи в материалах конференции}.

\bibitem{gupta2015}
Gupta S., Agrawal A., Gopalakrishnan K., Narayanan P.
Deep Learning with Limited Numerical Precision //
Proceedings of the 32nd International Conference on Machine Learning.
Proceedings of Machine Learning Research. 2015. Vol.~37. P.~1737--1746.
\href{https://proceedings.mlr.press/v37/gupta15.html}{Страница статьи в материалах конференции}.

\bibitem{flexpoint2017}
Köster U., Webb T.~J., Wang X. et al.
Flexpoint: An Adaptive Numerical Format for Efficient Training of Deep
Neural Networks // Advances in Neural Information Processing Systems.
2017. Vol.~30.
\href{https://arxiv.org/abs/1711.02213v2}{arXiv:1711.02213v2}.

\bibitem{fp82022}
Micikevicius P., Stosic D., Burgess N. et al.
FP8 Formats for Deep Learning. 2022.
\href{https://arxiv.org/abs/2209.05433v2}{arXiv:2209.05433v2}.

\bibitem{mx2023}
Darvish Rouhani B., Zhao R., More A. et al.
Microscaling Data Formats for Deep Learning. 2023.
\href{https://arxiv.org/abs/2310.10537v3}{arXiv:2310.10537v3}.

\bibitem{han2016eie}
Han S., Liu X., Mao H. et al.
EIE: Efficient Inference Engine on Compressed Deep Neural Network //
Proceedings of ISCA. 2016.
\href{https://arxiv.org/abs/1602.01528v2}{arXiv:1602.01528v2}.

\bibitem{teerapittayanon2016branchynet}
Teerapittayanon S., McDanel B., Kung H.~T.
BranchyNet: Fast Inference via Early Exiting from Deep Neural Networks //
International Conference on Pattern Recognition. 2016.
\href{https://doi.org/10.1109/ICPR.2016.7900006}{DOI: 10.1109/ICPR.2016.7900006}.

\bibitem{pragmatic2016}
Albericio J., Judd P., Delmas A., Sharify S., Moshovos A.
Bit-Pragmatic Deep Neural Network Computing. 2016.
\href{https://arxiv.org/abs/1610.06920}{arXiv:1610.06920}.

\bibitem{stripes2016}
Judd P. et al.
Stripes: Bit-Serial Deep Neural Network Computing // MICRO. 2016.
\href{https://doi.org/10.1109/MICRO.2016.7783722}{DOI: 10.1109/MICRO.2016.7783722}.

\bibitem{dynamic2017}
Delmas A., Judd P., Sharify S., Moshovos A.
Dynamic Stripes: Exploiting the Dynamic Precision Requirements
of Activation Values in Neural Networks. 2017.
\href{https://arxiv.org/abs/1706.00504}{arXiv:1706.00504}.

\bibitem{bitfusion2018}
Sharma H., Park J., Suda N. et al.
Bit Fusion: Bit-Level Dynamically Composable Architecture for
Accelerating Deep Neural Networks // ISCA. 2018.
\href{https://arxiv.org/abs/1712.01507v2}{arXiv:1712.01507v2}.

\bibitem{vit2021}
Dosovitskiy A., Beyer L., Kolesnikov A. et al.
An Image is Worth 16x16 Words: Transformers for Image Recognition at Scale //
ICLR. 2021. Электронная публикация: 2020.
\href{https://arxiv.org/abs/2010.11929}{arXiv:2010.11929}.

\bibitem{ho2020}
Ho J., Jain A., Abbeel P.
Denoising Diffusion Probabilistic Models //
Advances in Neural Information Processing Systems. 2020. Vol.~33.
\href{https://arxiv.org/abs/2006.11239}{arXiv:2006.11239}.

\bibitem{dit2023}
Peebles W., Xie S.
Scalable Diffusion Models with Transformers // ICCV. 2023.
Электронная публикация: 2022.
\href{https://arxiv.org/abs/2212.09748}{arXiv:2212.09748}.

\bibitem{mlperf2020}
Reddi V.~J., Cheng C., Kanter D. et al.
MLPerf Inference Benchmark // ISCA. 2020. Электронная публикация: 2019.
\href{https://arxiv.org/abs/1911.02549v2}{arXiv:1911.02549v2}.

\end{thebibliography}
\endgroup
